\section{Experiments}
\label{sec:experiments}
\vspace{-5.0pt}
In this section, we evaluate typicality bias across the modern vision stack and study its mitigation. We first provide our experimental setup in Sec.~\ref{sec:experimental_setup}, including datasets, evaluation metrics, models, and augmentation suites. We then survey five model families on the COCO NPMI splits and the synthetic STRANGE-Bench in Sec.~\ref{sec:results_diagnostic}, establishing that the bias persists across the stack and does not consistently shrink with scale. We finally run the systematic mitigation study on three representative families in Sec.~\ref{sec:results_mitigation}, showing that AAS reduces the typicality gap while preserving overall performance. Sec.~\ref{sec:results_human} checks with human raters that the gaps come from the models and not from the rendered scenes.

\subsection{Experimental Setup}
\label{sec:experimental_setup}

\noindent\textbf{Datasets and metrics.}
We evaluate on COCO \texttt{val}, Objects365v2 \texttt{val}~\cite{shao2019objects365} (mapped to COCO's 80-class label space for cross-domain transfer), and the synthetic STRANGE-Bench of Sec.~\ref{sec:strange}. On COCO and O365, we report overall performance as a sanity check (typicality-bias mitigation should not regress overall accuracy) and the typicality gap $\Delta^{\text{T-A}} = M^{\text{T}} - M^{\text{A}}$ on the per-anchor NPMI \emph{quintile} ($\alpha = 0.2$) splits of Sec.~\ref{sec:diagnostic}, where $M$ is AP for detection and micro-averaged F1 at a $0.5$ threshold (F1$_{\text{micro}}$) for DINOv3 classification. We compare rows within a family only, since the metrics (AP, F1) and their operating ranges differ. On STRANGE-Bench, we report the per-tier average $T_{1\text{-}5}^{\text{T}}, T_{1\text{-}5}^{\text{A}}$ across tiers $T_1$ to $T_5$ ($k = 0 \dots 4$) and the single-point $T_{\text{never}}$ score, the most extreme atypical case by construction (see Appendix~\ref{app:additional_results} for more details). The main paper uses $11$ of the $17$ sub-tiers ($T_1$ to $T_5$ in both arms and $T_{\text{never}}$); Tab.~\ref{tab:codino_extra_tiers} and Fig.~\ref{fig:strange_diagnostic_t6} report the others.

\begin{wrapfigure}{r}{0.50\linewidth}
    \vspace{-10pt}
    \centering
    \includegraphics[width=\linewidth]{static/images/strange_bench_per_tier_deltas_v2s.pdf}
    \captionsetup{skip=3pt}
    \caption{\textbf{Typicality bias is universal across modern vision models.} Per-tier typicality gap $\Delta = M^{\text{T}} - M^{\text{A}}$ on STRANGE-Bench for six models. The gap grows at every tier for the DINOv3 ViT-B linear probe (up to $\Delta = 37.2$ at T5), CLIP and OpenCLIP. Co-DINO R-50 rises to about $8.0$ by T3 and stays there. The LLaVA-OV gaps (F1 macro) are smaller.}
    \label{fig:strange_diagnostic}
    \vspace{-10pt}
\end{wrapfigure}

\noindent\textbf{Models.} We evaluate typicality bias on a representative model from each of five families spanning the modern vision stack: \emph{specialist detection} (Co-DINO~\cite{zong2023detrs} R-50, Swin-L), \emph{open-vocabulary detection} (Grounding DINO~\cite{liu2024grounding} Swin-B), \emph{self-supervised classification} (DINOv3~\cite{simeoni2025dinov3} ViT-B, both linear-probe and fine-tuned), \emph{vision--language models} (CLIP~\cite{radford2021learning} ViT-B/16, OpenCLIP~\cite{cherti2023reproducible} ViT-B/16 LAION-2B), and \emph{multimodal LLMs} (LLaVA-OneVision-7B~\cite{li2024llava}). 
For the diagnostic survey, we use the public pretrained checkpoints officially released by the authors. These families span the modern vision stack along two axes: \emph{supervision regime} (fully supervised, language-grounded, self-supervised, contrastive image--text, and instruction-tuned) and \emph{pretraining data scale} (from $\sim 100$K images for COCO-trained specialists to $\sim 2$ billion image--text pairs for OpenCLIP and LLaVA-OneVision). Any persistence of bias across this set is therefore unlikely to be an artifact of a single paradigm or scale.

For the systematic mitigation study, we select one representative from each of the three task settings, spanning supervision regime, architecture, and task: \textbf{Co-DINO R-50} (specialist closed-vocabulary detection), \textbf{Grounding DINO Swin-B} (open-vocabulary, text-conditioned), and \textbf{DINOv3 ViT-B/16 fine-tuned} (SSL foundational backbone fine-tuned for classification). All representative models are trained on COCO~\texttt{train2017} (initialized from the public pretrained checkpoints) under the augmentation variants below. Full training recipes, checkpoint sources, hyperparameters, and compute budgets are provided in Appendix~\ref{app:experimental-details}. We leave vision--language models and multimodal LLMs out of the mitigation study, since their co-occurrence prior is set during web-scale pretraining.

\noindent\textbf{Augmentation suite.}
We compare AAS against a representative suite of seven augmentations
spanning two axes, \emph{static}/\emph{dynamic} and \emph{non-generative}/\emph{generative}: \textbf{Color jitter}, \textbf{MixUp}~\cite{zhang2017mixup}, \textbf{Mosaic}~\cite{bochkovskiy2020yolov4}, \textbf{YOLOX}~\cite{ge2021yolox} (MixUp + Mosaic), \textbf{SCP}~\cite{ghiasi2021simple}, \textbf{XPaste}~\cite{zhao2023x}, and \textbf{InstAug}~\cite{kupyn2024dataset}. We additionally report a \textbf{Std.} baseline using only standard augmentation (flip, crop, resize). None of the existing augmentations are designed to break co-occurrence shortcuts; whether they incidentally affect typicality bias is an empirical question we answer below. AAS hyperparameters are selected once on a $25\%$ COCO subset and held fixed across architectures and tasks; per-augmentation hyperparameters are detailed in Appendix~\ref{app:exp-recipes}.

\subsection{Typicality Bias Persists Across the Modern Vision Stack}
\label{sec:results_diagnostic}
\vspace{-3.0pt}
To isolate the bias contained in the \emph{pretrained representation}
from any introduced by downstream task adaptation, we evaluate each family with the pretrained representation \emph{frozen}: zero-shot for Grounding DINO, CLIP, and OpenCLIP; linear probing for DINOv3; and
the released checkpoint for Co-DINO. We
survey the five model families on the COCO quintile split
(Fig.~\ref{fig:teaser}, right) and on the STRANGE-Bench per-tier breakdown (Fig.~\ref{fig:strange_diagnostic}).
The COCO quintile gap $\Delta^{\text{T-A}}$ is non-zero for every family we examine, ranging from \textbf{5.4}~p.p.\@ (Co-DINO Swin-L $1\times$) to \textbf{11.0}~p.p.\@ (OpenCLIP ViT-B/16); concretely, specialist detection shows $\Delta^{\text{T-A}}_{\text{AP}} = $ \textbf{9.5}~p.p.\@
(Co-DINO R-50), \textbf{5.4}~p.p.\@ (Swin-L $1\times$), and \textbf{5.6}~p.p.\@
(Swin-L $3\times$), open-vocabulary detection
$\Delta^{\text{T-A}}_{\text{AP}} = $ \textbf{6.3}~p.p.\@ (Grounding DINO
Swin-B), self-supervised classification
$\Delta^{\text{T-A}}_{\text{F1}}$ $= $ \textbf{8.6}~p.p.\@ (DINOv3 ViT-B
linear probe), and vision--language models
$\Delta^{\text{T-A}}_{\text{mAP}} = $ \textbf{9.5}~p.p.\@ (CLIP ViT-B/16)
and \textbf{11.0}~p.p.\@ (OpenCLIP ViT-B/16 LAION-2B). The Co-DINO Swin-L checkpoint, pretrained on a multi-dataset blend~\cite{zong2023detrs}, also exhibits $\Delta^{\text{T-A}}_{\text{AP}} = $ \textbf{5.7}~p.p.\@ on COCO val, indicating the bias is a property of the training distribution rather than dataset overfitting. From CLIP (400M image-text pairs) to OpenCLIP (LAION-2B), more pretraining data \emph{widens} rather than shrinks the gap (\textbf{9.5} $\rightarrow$ \textbf{11.0}~p.p.\@); scale alone is not a solution (App.~\ref{app:scale_ladder} separates data from capacity).
STRANGE-Bench amplifies this picture (Fig.~\ref{fig:strange_diagnostic}): the gap grows with contextual load $k$ at every tier for the DINOv3 linear probe, CLIP and OpenCLIP, reaching $\Delta = 37.2$ for the DINOv3 linear probe at the highest tier; Co-DINO R-50 rises to about $8.0$ by T3 and then stays flat, and the LLaVA-OV gaps (F1 macro) are smaller and do not grow with $k$. Per-tier breakdowns are provided in App.~\ref{app:additional_results}.

\begin{finding}
\textbf{Diagnosis.} Typicality bias persists across all five model families and four orders of magnitude of pretraining data ($\sim$100K $\rightarrow$ $\sim$2B images); scale and architectural diversity alone do not solve it.
\end{finding}

\subsection{Anchor-Atypicality Sampling Closes the Typicality Gap}
\label{sec:results_mitigation}

We evaluate the augmentation suite on the three representative families. An augmentation mitigates typicality bias when it lifts the atypical side while the typical side and overall accuracy hold; a smaller gap that comes from lowering both sides, as with Mosaic on Co-DINO (Tab.~\ref{tab:codino_main}), is not desirable. On STRANGE-Bench, the gap of the Std.\ baselines grows with the number of partners because the atypical side falls while the typical side changes little: from T1 to T5, Co-DINO goes from $90.8$ to $85.5$ AP on typical scenes and from $87.9$ to $77.3$ on atypical ones, DINOv3 from $91.3$ to $88.6$ F1 and from $92.0$ to $48.0$ (per-tier tables in App.~\ref{app:codino_full} and~\ref{app:dinov3_full}). Across all three, AAS preserves overall accuracy, closes
the typicality gap on the majority of (model, eval-split) cells, and
improves extreme-atypical performance the most.

\begin{table*}[t]
\centering
\small
\caption{\textbf{Co-DINO ResNet-50: typicality bias under augmentation
on COCO and STRANGE-Bench.} SCP + AAS \emph{(ours, blue)} has the highest COCO AP, and XPaste + AAS \emph{(ours, blue)} achieves
the smallest STRANGE-Bench tier gap, and the
highest never-co-occurring tier score; on COCO, Mosaic and YOLOX
attain the smallest typicality gaps, but do so by uniformly degrading both splits. Full-annotation bbox mAP (\%); $\Delta$ is typical minus atypical, $T_{1\text{-}5}$ the mean over tiers T1 to T5 of STRANGE-Bench. Best per column in \textbf{bold}, second-best \underline{underlined}.}
\vspace{-2pt}
\label{tab:codino_main}
\begin{tabular}{l cccc cccc}
\toprule
& \multicolumn{4}{c}{\textbf{COCO}} & \multicolumn{4}{c}{\textbf{STRANGE-Bench}} \\
\cmidrule(lr){2-5} \cmidrule(lr){6-9}
Augmentation & $\text{AP}^{\text{val}}$ $\uparrow$ & $\text{AP}^{\text{T}}$ $\uparrow$ & $\text{AP}^{\text{A}}$ $\uparrow$ & $\Delta^{\text{T-A}}$ $\downarrow$ & $T_{1\text{-}5}^{\text{T}}$ $\uparrow$ & $T_{1\text{-}5}^{\text{A}}$ $\uparrow$ & $\Delta_{1\text{-}5}$ $\downarrow$ & $T_{\text{never}}$ $\uparrow$ \\
\midrule
No aug & 43.5 & 44.2 & 35.5 & 8.7 & 83.4 & 76.4 & 7.0 & 76.3 \\
\midrule
\multicolumn{9}{l}{\emph{Non-generative}} \\
Std. & 51.8 & \underline{53.2} & 44.0 & 9.2 & 88.2 & 81.9 & 6.3 & 83.3 \\
Color jitter & 47.6 & 49.1 & 40.1 & 9.0 & 86.6 & 79.4 & 7.2 & 77.0 \\
MixUp & 46.2 & 47.2 & 39.1 & 8.1 & 84.0 & 77.0 & 7.0 & 78.0 \\
Mosaic & 45.8 & 46.1 & 39.7 & \textbf{6.4} & 72.2 & 68.0 & 4.2 & 69.3 \\
YOLOX & 46.7 & 47.0 & 40.3 & \underline{6.7} & 83.2 & 78.8 & 4.3 & 79.4 \\
SCP & \underline{52.1} & \underline{53.2} & 44.2 & 9.0 & 90.2 & 86.3 & 3.9 & 87.5 \\
\midrule
\multicolumn{9}{l}{\emph{Generative}} \\
InstAug & 51.7 & 52.6 & \underline{44.3} & 8.3 & 87.8 & 82.6 & 5.2 & 81.9 \\
XPaste & 50.9 & 52.0 & 43.0 & 9.0 & 90.7 & 89.0 & \underline{1.7} & \underline{90.4} \\
\midrule
\rowcolor{oursrow} SCP + AAS \emph{(ours)} & \textbf{52.9} & \textbf{54.0} & \textbf{45.5} & 8.5 & \underline{91.4} & \underline{89.2} & 2.1 & 90.1 \\
\rowcolor{oursrow} \textbf{XPaste + AAS \emph{(ours)}} & 51.4 & 52.1 & 43.3 & 8.8 & \textbf{91.7} & \textbf{90.7} & \textbf{1.0} & \textbf{92.1} \\
\bottomrule
\end{tabular}

\vspace{-8pt}
\end{table*}


\noindent\textbf{Supervised detection.}
On Co-DINO R-50 (Tab.~\ref{tab:codino_main}), SCP + AAS achieves the highest AP on COCO (overall: \textbf{52.9}, typical: \textbf{54.0}, atypical: \textbf{45.5}; $+1.1$~p.p.\@, $+0.8$~p.p.\@ and $+1.5$~p.p.\@ over Std.) and XPaste + AAS on STRANGE-Bench (typical: \textbf{91.7}, atypical: \textbf{90.7}, $T_{\text{never}}$: \textbf{92.1}; $+3.5$~p.p.\@, $+8.8$~p.p.\@ and $+8.8$~p.p.\@ over Std.). AAS lowers the COCO quintile gap of the baseline it starts from on
both sides (SCP \textbf{9.0}~$\rightarrow$ \textbf{8.5}, XPaste \textbf{9.0}~$\rightarrow$ \textbf{8.8}) and raises overall AP
($52.1 \rightarrow 52.9$ and $50.9 \rightarrow 51.4$). On STRANGE-Bench the change is larger with AAS:
$\Delta_{1\text{-}5}$ falls from \textbf{3.9} to \textbf{2.1} ($-1.8$~p.p.\@) for SCP and from \textbf{1.7} to \textbf{1.0} ($-0.7$~p.p.\@) for XPaste, and
$T_{\text{never}}$ rises from \textbf{87.5} to \textbf{90.1} ($+2.6$~p.p.\@) and from \textbf{90.4} to \textbf{92.1} ($+1.7$~p.p.\@).
Mosaic and YOLOX reach the smallest COCO gaps ($6.4$ and $6.7$) by lowering AP on both splits
($-3.7$ to $-7.1$~p.p.\@ vs.\ Std.). AAS closes the
gap by lifting atypical performance above its Std.\ baseline. Per-tier
breakdowns and results on the cross-domain Objects365 split are provided in Appendix~\ref{app:codino_full}.

\noindent\textbf{Open-vocabulary detection.}
Table~\ref{tab:grounding_dino_main} reports Grounding DINO Swin-B
under five configurations. Three observations stand out. First,
the zero-shot baseline already exhibits a substantial COCO quintile
gap ($\Delta^{\text{T-A}} = 6.3$~p.p.\@); fine-tuning on COCO with paste augmentation raises COCO val AP by $1.6$ to $2.0$~p.p.\@ but keeps the gap at $5.6$ to $6.2$. Second, AAS does not lower COCO val AP relative to the row it starts from ($+0.1$~p.p.\@ for SCP, $+0.2$~p.p.\@ for XPaste), with XPaste + AAS the highest (\textbf{58.9}), and keeps the COCO gap within $0.4$~p.p.\@ of that row ($6.2 \rightarrow 6.2$, $5.6 \rightarrow 6.0$). Third,
on STRANGE-Bench, AAS lifts atypical AP over the row it starts from ($+1.3$~p.p.\@ for SCP, $+1.1$~p.p.\@ for XPaste) and $T_{\text{never}}$ ($+1.2$~p.p.\@, $+0.6$~p.p.\@), and XPaste + AAS attains the highest typical (\textbf{95.0}), atypical (\textbf{94.7}) and $T_{\text{never}}$ (\textbf{95.2}) scores ($+1.7$~p.p.\@, $+3.8$~p.p.\@ and $+2.8$~p.p.\@ over zero-shot) and, with XPaste, the smallest mean tier gap ($\Delta_{1\text{-}5} = $ \textbf{0.3}~p.p.\@; zero-shot: $2.4$~p.p.\@). 

\begin{table*}[h]
\centering
\small
\caption{\textbf{Grounding DINO Swin-B: typicality bias under
augmentation on COCO and STRANGE-Bench.} XPaste + AAS \emph{(ours,
blue)} attains the highest COCO val AP and the highest
STRANGE-Bench scores.}
\label{tab:grounding_dino_main}
\vspace{-2pt}
\begin{tabular}{l cccc cccc}
\toprule
& \multicolumn{4}{c}{\textbf{COCO}} & \multicolumn{4}{c}{\textbf{STRANGE-Bench}} \\
\cmidrule(lr){2-5} \cmidrule(lr){6-9}
Augmentation & $\text{AP}^{\text{val}}$ $\uparrow$ & $\text{AP}^{\text{T}}$ $\uparrow$ & $\text{AP}^{\text{A}}$ $\uparrow$ & $\Delta^{\text{T-A}}$ $\downarrow$ & $T_{1\text{-}5}^{\text{T}}$ $\uparrow$ & $T_{1\text{-}5}^{\text{A}}$ $\uparrow$ & $\Delta_{1\text{-}5}$ $\downarrow$ & $T_{\text{never}}$ $\uparrow$ \\
\midrule
Zero Shot & 56.9 & 57.7 & 51.4 & 6.3 & 93.3 & 90.9 & 2.4 & 92.4 \\
\midrule
\multicolumn{9}{l}{\emph{Non-generative}} \\
SCP & 58.5 & 59.6 & 53.4 & 6.2 & 93.3 & 92.4 & 0.9 & 93.3 \\
\midrule
\multicolumn{9}{l}{\emph{Generative}} \\
XPaste & \underline{58.7} & \underline{59.7} & \textbf{54.1} & \textbf{5.6} & 93.9 & 93.6 & \textbf{0.3} & \underline{94.6} \\
\midrule
\rowcolor{oursrow} SCP + AAS \emph{(ours)} & 58.6 & \textbf{59.8} & 53.6 & 6.2 & \underline{94.3} & \underline{93.7} & \underline{0.6} & 94.5 \\
\rowcolor{oursrow} \textbf{XPaste + AAS \emph{(ours)}} & \textbf{58.9} & \underline{59.7} & \underline{53.7} & \underline{6.0} & \textbf{95.0} & \textbf{94.7} & \textbf{0.3} & \textbf{95.2} \\
\bottomrule
\end{tabular}

\vspace{-3.0pt}
\end{table*}


\begin{table*}[t]
\centering
\small
\caption{\textbf{DINOv3 ViT-B/16: typicality bias under augmentation
on COCO and STRANGE-Bench.} XPaste + AAS \emph{(ours, blue)} achieves
the highest scores on both datasets and the smallest STRANGE-Bench typicality gap, with the largest gains over Std.\ on atypical scenes ($+2.5$~p.p.\@ on COCO, $+28.0$~p.p.\@ on STRANGE-Bench and $+33.5$~p.p.\@ on $T_{\text{never}}$); its COCO gap ($3.9$) is second to SCP and YOLOX ($3.7$). F1$_{\text{micro}}$ (\%), columns as in Tab.~\ref{tab:codino_main}.}
\label{tab:dinov3_main}
\vspace{-2pt}
\begin{tabular}{l cccc cccc}
\toprule
& \multicolumn{4}{c}{\textbf{COCO}} & \multicolumn{4}{c}{\textbf{STRANGE-Bench}} \\
\cmidrule(lr){2-5} \cmidrule(lr){6-9}
Augmentation & $\text{F1}^{\text{val}}$ $\uparrow$ & $\text{F1}^{\text{T}}$ $\uparrow$ & $\text{F1}^{\text{A}}$ $\uparrow$ & $\Delta^{\text{T-A}}$ $\downarrow$ & $T_{1\text{-}5}^{\text{T}}$ $\uparrow$ & $T_{1\text{-}5}^{\text{A}}$ $\uparrow$ & $\Delta_{1\text{-}5}$ $\downarrow$ & $T_{\text{never}}$ $\uparrow$ \\
\midrule
Std.\ \emph{(linear probe)} & 74.6 & 71.8 & 63.2 & 8.6 & 63.5 & 39.6 & 23.9 & 38.9 \\
\midrule
No aug & 78.9 & 77.8 & 70.1 & 7.7 & 85.5 & 66.3 & 19.2 & 61.5 \\
\midrule
\multicolumn{9}{l}{\emph{Non-generative}} \\
Std. & 81.9 & 79.8 & 74.3 & 5.5 & 89.7 & 67.5 & 22.2 & 62.6 \\
Color jitter & 81.8 & 79.7 & 75.1 & 4.6 & 89.2 & 66.7 & 22.5 & 62.7 \\
MixUp & \underline{82.2} & 79.6 & 75.4 & 4.2 & 89.9 & 72.4 & 17.5 & 69.9 \\
Mosaic & 69.6 & 71.0 & 67.0 & 4.0 & 74.7 & 68.0 & 6.7 & 70.6 \\
YOLOX & 70.0 & 71.2 & 67.4 & \textbf{3.7} & 75.1 & 69.1 & 6.0 & 70.7 \\
SCP & 81.6 & 78.8 & 75.1 & \textbf{3.7} & 91.4 & 87.1 & 4.3 & 86.5 \\
\midrule
\multicolumn{9}{l}{\emph{Generative}} \\
InstAug & 81.9 & 79.5 & 74.4 & 5.2 & 89.5 & 67.7 & 21.8 & 63.3 \\
XPaste & 82.1 & \underline{80.2} & \underline{76.1} & 4.1 & \underline{93.8} & \underline{94.4} & \underline{$-0.7$} & \underline{95.0} \\
\midrule
\rowcolor{oursrow} SCP + AAS \emph{(ours)} & 81.7 & 79.0 & 75.0 & 4.0 & 92.5 & 91.2 & 1.2 & 89.5 \\
\rowcolor{oursrow} \textbf{XPaste + AAS \emph{(ours)}} & \textbf{82.5} & \textbf{80.6} & \textbf{76.8} & \underline{3.9} & \textbf{94.6} & \textbf{95.5} & \textbf{\boldmath $-0.8$} & \textbf{96.1} \\
\bottomrule
\end{tabular}

\vspace{-10pt}
\end{table*}


\noindent\textbf{SSL classification.}
Table~\ref{tab:dinov3_main} reports the full augmentation suite on
DINOv3 ViT-B/16 fine-tuned, with the linear probe shown in the top
row for reference. AAS yields the largest mitigation effect of the
three families: XPaste+AAS reduces the COCO quintile gap from
$\Delta^{\text{T-A}} = $ \textbf{5.5} (Std.) to \textbf{3.9} ($-1.6$
p.p.\@, a 29\% relative reduction), while improving overall F1
($81.9 \rightarrow 82.5$, $+0.6$ p.p.\@) and lifting the atypical-split
score by $+2.5$ p.p.\@ ($74.3 \rightarrow 76.8$). On STRANGE-Bench, the
gap is essentially eliminated: $\Delta_{1\text{-}5}$ collapses from
\textbf{22.2} to \textbf{$-0.8$} ($-23.0$ p.p.\@), and $T_{\text{never}}$
rises from $62.6$ to \textbf{96.1} ($+33.5$ p.p.\@). The non-paste variants
either leave the gap large (Color jitter, MixUp, InstAug) or reduce
it primarily by uniform degradation (Mosaic, YOLOX collapse overall
F1 by $\sim 12$ p.p.\@), while SCP and XPaste already narrow it (COCO $3.7$ and $4.1$, $\Delta_{1\text{-}5}$ $4.3$ and $-0.7$).
The situation differs for linear probing on DINOv3 ViT-B/16 with the backbone frozen, where augmentations operate only on the head input and their effect is muted; we report and discuss these results in Appendix~\ref{app:dinov3_lp}.

\begin{finding}
\textbf{Mitigation.} AAS closes the typicality gap by up to
$\mathbf{1.6}$ p.p.\@ on COCO splits and $\mathbf{23.0}$ p.p.\@ on
STRANGE-Bench, while preserving overall performance: the
intervention is asymmetric, lifting atypical-context performance
without degrading typical-context performance.
\end{finding}

\subsection{The Typicality Gap Is Not a Rendering Artifact}
\label{sec:results_human}
\vspace{-3.0pt}
\begin{wraptable}{r}{0.33\linewidth}
\vspace{-11pt}
\centering\footnotesize
\setlength{\tabcolsep}{3pt}
\captionsetup{font=small}
\caption{\textbf{Free listing}: anchor recall (\%); models: DINOv3 fine-tuned.}
\label{tab:human_listing}
\begin{tabular}{lccc}
\toprule
Rater & $T_1$ & $T_{\text{never}}$ & $\Delta$ \\
\midrule
Humans & $98.3$ & $97.3$ & $-1.0$ \\
\midrule
Std. & $98.0$ & $67.0$ & $-31.0$ \\
XPaste & $100.0$ & $96.0$ & $-4.0$ \\
\rowcolor{oursrow} XPaste + AAS & $100.0$ & $97.0$ & $-3.0$ \\
\bottomrule
\end{tabular}

\vspace{-10pt}
\end{wraptable}
Three blinded user studies test whether the gaps on STRANGE-Bench come from the renders or from the models (App.~\ref{app:strange:user-study}); they ran on an initial version of the benchmark with one render per scene and a lighter audit. In \emph{free listing}, raters named every category they could identify in $100$ $T_1$ and $100$ $T_{\text{never}}$ scenes shown without boxes. A rendering defect would hide the anchor from humans and models alike, yet human anchor recall drops by $1.0$~p.p.\@ and that of the DINOv3 Std.\ model by $31.0$~p.p.\@ (Tab.~\ref{tab:human_listing}). In a \emph{quality audit} of $330$ scenes, typical and atypical scenes at the same $k$ differ in pass rate by $2.2$~p.p.\@ on average over four criteria, with no systematic direction.
